\documentclass{article}
\usepackage{tikz}
\usetikzlibrary{shadings} % LATEX and plain TEX when using TikZ
\begin{document}
\tikz \draw[top color=red] (0,0) rectangle (2,1);

\tikz \draw[top color=white,bottom color=black,middle color=red]
(0,0) rectangle (2,1);

\begin{tikzpicture}
\shade[ball color=white] (0,0) circle (2ex);
\shade[ball color=red] (1,0) circle (2ex);
\shade[ball color=black] (2,0) circle (2ex);
\end{tikzpicture}

% \tikz
% \shade[upper left=red,upper right=green,
% lower left=blue,lower right=yellow]
% (0,0) rectangle (3,2);

\tikz \shade[shading=color wheel] (0,0) circle (1.5);

\tikz \shade[shading=color wheel] [even odd rule]
(0,0) circle (1.5)
(0,0) circle (1);

\tikz \shade[shading=color wheel black center] (0,0) circle (1.5);

\tikz \shade[shading=color wheel white center] (0,0) circle (1.5);

\tikz \shade[shading=Mandelbrot set] (0,0) rectangle (2,2);

\tikz \draw[inner color=red] (0,0) rectangle (2,1);

\tikz \draw[outer color=red,inner color=white]
(0,0) rectangle (2,1);

\end{document}
